\documentclass[aspectratio=169,11pt]{beamer}

% =========================================================
% AUTOMATISMES -- STATISTIQUES 4e
% Une série par diapositive.
% Mise en page simple et lisible en classe.
% =========================================================

\usepackage[T1]{fontenc}
\usepackage[utf8]{inputenc}
\usepackage[french]{babel}
\usepackage{lmodern}
\usepackage{amsmath,amssymb,array}

% Palette sobre
\definecolor{TitreC}{RGB}{39,55,80}
\definecolor{SavoirC}{RGB}{50,105,170}
\definecolor{NeutreC}{RGB}{80,88,100}

\setbeamertemplate{navigation symbols}{}
\setbeamertemplate{footline}{%
  \hfill
  {\scriptsize\color{NeutreC}\insertframenumber}
  \hspace{0.4cm}\vspace{0.2cm}
}

\setbeamercolor{background canvas}{bg=white}
\setbeamercolor{normal text}{fg=TitreC,bg=white}
\setbeamercolor{frametitle}{fg=TitreC,bg=SavoirC!10}

\setbeamerfont{normal text}{size=\large}
\setbeamerfont{frametitle}{size=\Large,series=\bfseries}

\setbeamersize{text margin left=0.8cm,text margin right=0.8cm}

\newcommand{\serie}[2]{%
\begin{frame}[t]{#1}
\vspace{0.25cm}
#2
\end{frame}
}

\begin{document}

% =========================================================
% TITRE
% =========================================================
\begin{frame}[plain]
\vfill
\centering
{\Huge\bfseries\color{TitreC} Automatismes -- Statistiques\par}
\vspace{0.5cm}
{\Large\color{SavoirC}4\ieme{} -- Séries courtes}
\vfill
\end{frame}

% =========================================================
% SÉRIE 1
% =========================================================
\serie{Série 1 -- Lire des données brutes}{

On considère la série :

\[
2\quad3\quad1\quad2\quad4\quad2\quad0\quad3\quad2\quad1.
\]

\begin{enumerate}
  \item Quel est l'effectif total ?
  \item Quel est l'effectif de la valeur \(2\) ?
  \item Quelle est la valeur la plus fréquente ?
\end{enumerate}
}

% =========================================================
% SÉRIE 2
% =========================================================
\serie{Série 2 -- Fréquences}{

Dans une classe de \(25\) élèves, \(10\) viennent à pied.

\begin{enumerate}
  \item Écris la fréquence sous forme d'une fraction.
  \item Donne son écriture décimale.
  \item Donne le pourcentage correspondant.
\end{enumerate}
}

% =========================================================
% SÉRIE 3
% =========================================================
\serie{Série 3 -- Contrôler un tableau}{

Un tableau d'effectifs contient :

\[
4\qquad7\qquad5\qquad3\qquad6.
\]

\begin{enumerate}
  \item Calcule l'effectif total.
  \item Les fréquences suivantes peuvent-elles convenir ?
\end{enumerate}

\[
16\%\qquad28\%\qquad20\%\qquad12\%\qquad24\%.
\]
}

% =========================================================
% SÉRIE 4
% =========================================================
\serie{Série 4 -- Moyenne simple}{

Calcule puis interprète la moyenne de :

\[
8\qquad10\qquad12\qquad14.
\]
}

% =========================================================
% SÉRIE 5
% =========================================================
\serie{Série 5 -- Retrouver une valeur}{

La moyenne de cinq valeurs est \(12\).

Quatre de ces valeurs sont :

\[
9\qquad11\qquad13\qquad15.
\]

\medskip

Quelle est la cinquième valeur ?
}

% =========================================================
% SÉRIE 6
% =========================================================
\serie{Série 6 -- Moyenne avec effectifs}{

\[
\begin{array}{c|ccc}
\text{Valeur} & 5 & 10 & 15\\
\hline
\text{Effectif} & 2 & 3 & 1
\end{array}
\]

\medskip

Calcule la moyenne.
}

% =========================================================
% SÉRIE 7
% =========================================================
\serie{Série 7 -- Médiane, effectif impair}{

Détermine puis interprète la médiane :

\[
9\qquad4\qquad12\qquad7\qquad10.
\]
}

% =========================================================
% SÉRIE 8
% =========================================================
\serie{Série 8 -- Médiane, effectif pair}{

Détermine la médiane :

\[
3\qquad8\qquad5\qquad11\qquad6\qquad14.
\]
}

% =========================================================
% SÉRIE 9
% =========================================================
\serie{Série 9 -- Médiane à partir d'effectifs}{

\[
\begin{array}{c|ccccc}
\text{Valeur} & 1 & 2 & 3 & 4 & 5\\
\hline
\text{Effectif} & 2 & 4 & 5 & 3 & 1
\end{array}
\]

\begin{enumerate}
  \item Quel est l'effectif total ?
  \item Quelle position occupe la médiane ?
  \item Détermine la médiane.
\end{enumerate}
}

% =========================================================
% SÉRIE 10
% =========================================================
\serie{Série 10 -- Pourcentages et angles}{

Complète mentalement :

\[
50\%\longleftrightarrow\ldots^\circ
\qquad\qquad
25\%\longleftrightarrow\ldots^\circ
\]

\[
10\%\longleftrightarrow\ldots^\circ
\qquad\qquad
20\%\longleftrightarrow\ldots^\circ
\]
}

% =========================================================
% SÉRIE 11
% =========================================================
\serie{Série 11 -- Angle d'un secteur}{

Dans une enquête portant sur \(40\) personnes,
\(14\) choisissent la réponse A.

\medskip

Calcule l'angle du secteur correspondant.
}

% =========================================================
% SÉRIE 12
% =========================================================
\serie{Série 12 -- Moyenne ou médiane ?}{

Une entreprise a les salaires suivants :

\[
1500\quad1550\quad1600\quad1650\quad1700\quad9000.
\]

\medskip

Sans tout calculer, quel indicateur sera le plus influencé
par le salaire de \(9000\) euros ?
}

% =========================================================
% SÉRIE 13
% =========================================================
\serie{Série 13 -- Vrai ou faux ?}{

Justifie chaque réponse en une phrase.

\begin{enumerate}
  \item La somme des fréquences d'une série vaut \(100\).
  \item La médiane appartient toujours à la série.
  \item Une moyenne peut ne pas être une valeur observée.
  \item Dans un diagramme circulaire, la somme des angles vaut \(360^\circ\).
\end{enumerate}
}

% =========================================================
% SÉRIE 14
% =========================================================
\serie{Série 14 -- Contrôler une affirmation}{

Dans une série, la médiane vaut \(17\).

\medskip

Complète une phrase correcte commençant par :

\begin{center}
\Large\bfseries
« Au moins la moitié des valeurs... »
\end{center}
}

\end{document}
